\section{Oracoli e Meccanismi di
Verifica}\label{oracoli-e-meccanismi-di-verifica}

Per decidere se un team ha recuperato la \textbf{flag}, bisogna definire
un\textquotesingle{}\emph{\textbf{oracolo}}, che determini lo stato in
cui si trova il sistema target.

L\textquotesingle oracolo e\textquotesingle{} in grado di accedere al
sistema e di fare due azioni:

\begin{itemize}
\tightlist
\item
  puo\textquotesingle{} leggere la memoria fisica di ogni nodo
\item
  puo\textquotesingle{} eseguire comandi su qualsiasi nodo
\end{itemize}

Per ogni tipo di azione eseguita, l\textquotesingle oracolo determina lo
stato del sistema confrontando il risultato dell\textquotesingle azione
con il valore che l\textquotesingle oracolo si aspettava (ossia i valori
vengono validati).

Dunque all\textquotesingle iterno del processo di verifica, il sistema
si puo\textquotesingle{} definire \emph{patchato} se e soltanto se tutte
le azioni sono validate rispetto al valore che ci si aspettava in
output.

Di seguito sono discussi alcuni meccanismi di verifica che abbiamo
identificato.

\subsection{Verifica Statica}\label{verifica-statica}

La verifica statica prevede che l\textquotesingle oracolo possa accedere
alla memoria fisica, comparando le patch applicate nei file in memoria,
ad una patch considerata di riferimento per l\textquotesingle oracolo.

Questo tipo di verifica puo\textquotesingle{} essere implementato in
vari modi:

\begin{itemize}
\tightlist
\item
  \textbf{diff}: ad ogni file interessato corrisponde un file di
  riferimento per l\textquotesingle oracolo. Ogni file viene dunque
  confrontato, riga per riga, con il suo riferimento per determinare se
  il team ha applicato la patch correttamente.
\item
  \textbf{SAST}: Static Analysis Security Testing. Prevede
  l\textquotesingle accesso al codice sorgente non compilato in modo da
  analizzare pattern di codice vulnerabili e fallace logiche.\footnote{\label{fn:ref_1}\url{https://aetic.theiaer.org/archive/v4/v4n3/p1.html}
    2. Currrent Status of SAST based on Checkmarx}
\end{itemize}

Tuttavia, SAST non e\textquotesingle{} un metodo di analisi efficace,
sebbene efficiente:

\begin{itemize}
\tightlist
\item
  possono essere identificati nel processo di verifica, vari falsi
  positivi (e falsi negativi)\footref{fn:ref_1}
\item
  non posso identificare staticamente tutti i tipi di errori di
  sicurezza \footref{fn:ref_1}
\item
  non posso identificare vulnerabilita all\textquotesingle interno dei
  binari compilati\footref{fn:ref_1}
\end{itemize}

Dunque l\textquotesingle analisi statica puo\textquotesingle{} in alcuni
casi risultare inesetta, e non coprire tutti i casi
d\textquotesingle uso per quanto riguarda la gestione di competizioni
CTF.

\subsection{Verifica Dinamica tramite sistema
esterno}\label{verifica-dinamica-tramite-sistema-esterno}

Questo meccanismo di verifica prevede l\textquotesingle uso di un
\textbf{orchestratore esterno} alle macchine vulnerabili, ossia
l\textquotesingle attaccante, che possa interagire con la vulnerabilita
utilizzando la rete locale. L\textquotesingle attaccante orchestra
l\textquotesingle esecuzione di uno o piu comandi, con
l\textquotesingle obiettivo di inviare il payload
all\textquotesingle interno di un pacchetto di rete.

Durante una competizione di questo tipo, tuttavia, i membri dei team
sono legittimati ad accedere alla macchina vulnerabile come utente
\passthrough{\lstinline!root!}: cosi da permettere
l\textquotesingle applicazione delle patch sui file di sistema. Questo
implica che i partecipanti dei team siano in grado di eseguire tool come
\passthrough{\lstinline!tcpdump!} e \passthrough{\lstinline!wireshark!}
per intercettare il traffico di rete e di conseguenza, di visualizzarne
il payload, ottenendo un suggerimento per la soluzione della challenge.

Inoltre, il payload inviato potrebbe essere tale da mandare in crash il
servizio target, causando una penalita alla percentuale di \textbf{SLA}
\emph{(Service Level Agreement)}, ed interferendo potenzialmente con il
lavoro del team sfidante (portando per esempio alla perdita di alcuni
file). Si vuole dunque trovare un sistema di verifica che esegua
completamente all\textquotesingle oscuro del team sfidante.

Il sistema di verifica tramite sistema esterno non copre nemmeno tutte
le vulnerabilita. Banalmente, non tutti i servizi vulnerabili sono
esposti in rete, ma sono accessibili solo in locale.

Riassumendo quindi, abbiamo identificato alcuni spunti per ideare un
meccanismo di verifica migliore:

\begin{itemize}
\tightlist
\item
  Il traffico di rete non dovrebbe essere letto dal team sfidante al
  fine di ottenere suggerimenti per la soluzione.
\item
  Se l\textquotesingle exploit manda in crash il servizio, questo non
  dovrebbe penalizzare il team, o piu in generale: il team sfidante non
  vede il servizio distruggersi.
\item
  Dobbiamo gestire anche vulnerabilita non esposte in rete.
\end{itemize}

\subsection{Verifica Dinamica tramite injection dei comandi di
verifica}\label{verifica-dinamica-tramite-injection-dei-comandi-di-verifica}

Questo meccanismo di verifica prevede l\textquotesingle uso di un
\textbf{orcehstratore esterno} al sistema vulnerabile, che ha la facolta
di eseguire comandi all\textquotesingle interno delle macchine
vulnerabili che compongono il sistema. Dunque
l\textquotesingle orchestratore esegue uno o piu comandi con il fine di
verificare che il sistema sia resistente alla vulnerabilita.

Considerando il caso in cui una macchina \passthrough{\lstinline!A!}
contenga un file binario vulnerabile ad un \emph{buffer overflow},
l\textquotesingle orchestratore potrebbe eseguire ad esempio il comando
\passthrough{\lstinline!/path/to/service\_logger payload!} e verificare
l\textquotesingle esito dell\textquotesingle esecuzione.

L\textquotesingle esecuzione di questi comandi e\textquotesingle{}
possibile per esempio in ambienti containerizzati, sfruttando utility
come \passthrough{\lstinline!docker exec!} o
\passthrough{\lstinline!incus exec!}, che permettono di iniettare
comandi all\textquotesingle interno dei container.

Tuttavia: comandi come \passthrough{\lstinline!pspy!}, (che ascoltano
attivamente le chiamate di sistema effettuate e gli eventi sul file
system in \passthrough{\lstinline!/proc!}) possono intercettare
l\textquotesingle esecuzione dei comandi
dell\textquotesingle orchestratore e di mostrarli al team, suggerendo al
team sfidante la soluzione. \footnote{\label{fn:ref_2}\url{https://github.com/DominicBreuker/pspy}}

Dunque, dobbiamo ideare un meccanismo di verifica che deve sia
completamente invisibile a questi meccanismi di intercettazione.

\subsection{Clonazione del sistema in una rete
separata}\label{clonazione-del-sistema-in-una-rete-separata}

In base a quanto visto prima, in BlueAgent faremo cosi:

\begin{enumerate}
\def\labelenumi{\arabic{enumi}.}
\tightlist
\item
  Si clonano tutte le macchine del sistema vulnerabile, in una rete
  separata, mantenendone la topologia originale.
\end{enumerate}

\begin{itemize}
\tightlist
\item
  \begin{itemize}
  \tightlist
  \item
    La rete di verifica per il team \(k\)-esimo sara\textquotesingle{}
    chiamata \(\text{verify-}k\)
  \end{itemize}
\end{itemize}

\begin{enumerate}
\def\labelenumi{\arabic{enumi}.}
\setcounter{enumi}{1}
\tightlist
\item
  Gli sfidanti non devono poter instradare i pacchetti da e verso la
  rete \(\text{verify-}k\).
\item
  Gli sfidanti non devono poter fare injection dei comandi nel sistema
  nella rete \(\text{verify-}k\).
\item
  Si aspetta che il sistema si operativo, facendo dei controlli di
  readyness.
\item
  Si esegue una batteria di test sul sistema clonato. Questo test
  e\textquotesingle{} costituita da una serie di comandi da iniettare
  nel sistema clonato.
\item
  Per ogni comando eseguito si comparano
  \passthrough{\lstinline!stdout!}, \passthrough{\lstinline!stderr!} ed
  \passthrough{\lstinline!exit\_code!} con il valore attesso.
\item
  Se il test ha avuto successo, allora BlueAgent risponde con lo stato
  \passthrough{\lstinline!PATCHED!}. (TODO: CONTROLLA CHE SIA CORRETTO!)
\end{enumerate}
